% Generated draft: requires mathematical and source review.
\documentclass[11pt]{article}
\usepackage[margin=25mm]{geometry}
\usepackage{fontspec}
\setmainfont{Latin Modern Roman}
\usepackage{amsmath,amssymb,mathtools,unicode-math}
\setmathfont{Latin Modern Math}
\usepackage{enumitem,needspace,xurl,hyperref}
\setlength{\parindent}{0pt}
\setlength{\parskip}{5pt}
\setlength{\emergencystretch}{3em}
\newcommand{\N}{\mathbb{N}}
\newcommand{\Z}{\mathbb{Z}}
\newcommand{\Q}{\mathbb{Q}}
\newcommand{\R}{\mathbb{R}}
\newcommand{\C}{\mathbb{C}}
\newcommand{\F}{\mathbb{F}}
\newcommand{\A}{\mathbb{A}}
\newcommand{\PP}{\mathbb{P}}
\newcommand{\E}{\mathbb{E}}
\newcommand{\eps}{\varepsilon}
\newcommand{\dd}{\,\mathrm{d}}
\begin{document}
\section*{UnsolvedMath ID: 30005035}
\textbf{Source problem code:} catalogue entry
\subsection{Definitions and mathematical statement}
The supplied clean statement concerns all finitely presented groups. The following gives a precise formulation under the integral cellular convention for homological filling; the supplied record does not specify the coefficient convention.

Let \(H\) be a group with a finite presentation \(\mathcal P=\langle A\mid R\rangle\). Let \(X_{\mathcal P}\) be the presentation complex and let \(\widetilde X_{\mathcal P}\) be its universal cover. For \(i=1,2\), let \(C_i(\widetilde X_{\mathcal P};\Z)\) denote the finitely supported integral cellular chains, with cellular boundary operator \(\partial_i\). For a chain \(c=\sum_\sigma a_\sigma\sigma\), define \(\lVert c\rVert_1=\sum_\sigma |a_\sigma|\). Write \(\N_0=\{0,1,2,\ldots\}\). For an integral cellular \(1\)-cycle \(z\), define
\[
\operatorname{Fill}_{\mathcal P}(z)
=\min\bigl\{\lVert b\rVert_1:b\in C_2(\widetilde X_{\mathcal P};\Z),\ \partial_2b=z\bigr\}.
\]
Such a filling exists because the universal cover is simply connected. Define the integral homological filling function by
\[
\operatorname{FA}_{\mathcal P}(n)
=\sup\bigl\{\operatorname{Fill}_{\mathcal P}(z):
 z\in\ker\partial_1,\ \lVert z\rVert_1\le n\bigr\},
\qquad n\in\N_0.
\]
Here \(\operatorname{FA}_H\) is interpreted through a finite presentation as above. A recursive bound means a total computable function \(f:\N_0\to\N_0\) satisfying \(\operatorname{FA}_{\mathcal P}(n)\le f(n)\) for every \(n\in\N_0\). The word problem for \(H\) is solvable if an algorithm halts on every finite word in \(A\cup A^{-1}\) and decides whether that word represents the identity of \(H\).

\emph{Question under this convention.} For every finitely presented group \(H\), does the existence of a finite presentation \(\mathcal P\) of \(H\) and a total computable function \(f:\N_0\to\N_0\) such that
\[
\forall n\in\N_0,\qquad \operatorname{FA}_{\mathcal P}(n)\le f(n)
\]
imply that \(H\) has a solvable word problem?

\emph{Unresolved source qualification.} This universal question formalizes the supplied clean statement. The supplied original statement asks only, ``Does \(H\) have solvable word problem?'', without defining that particular \(H\). Although it says that the report passage is retained verbatim, no identifying passage is included in the supplied record. The supplied literature assessment refers to an associated Leary group family whose criterion depends on a base group and on whether a parameter \(S\) is recursive; neither that family nor those parameters are defined here. Consequently, equivalence between the universal formulation above and the original source-specific question is unresolved.
\subsection{Short English statement}
Under the integral cellular interpretation of homological filling, the supplied clean question asks whether every finitely presented group with a computable upper bound on its homological filling function has a decidable word problem. The original extract leaves \(H\) unidentified, so its equivalence to this universal question requires review. The record labels the problem open; this is supplied curation, not an independently verified status.
\subsection{Sources}
\begin{itemize}
\item Local UnsolvedMath registry: supplied record with permanent ID \(30005035\), title \emph{Recursive Homological Filling and the Word Problem}, source code \texttt{catalogue entry}, in the Oberwolfach Reports Open Problems collection. The record gives proposed year 2022 and identifies no proposer.
\item Supplied report citation: \emph{Geometric Structures in Group Theory (2023)}, Oberwolfach Reports, pp.\ 517--576, DOI \texttt{10.4171/owr/2022/11}. \url{https://doi.org/10.4171/owr/2022/11}. The supplied background describes Oberwolfach Reports as documenting mathematical workshops and their open-problem discussions. No problem-specific page or complete contextual passage is supplied.
\item Supplied literature source: \emph{Homological Dehn functions of groups of type FP2}. \url{https://arxiv.org/abs/2012.00730}. No authors or theorem identifier are supplied in this record.
\item Supplied literature source: \emph{Oberwolfach report}. \url{https://doi.org/10.4171/owr/2022/11}.
\item Literature assessment checked at \texttt{2026-08-22}: this exact date records supplied dated curation, not independent verification. The supplied status is \texttt{open}. The assessment says that the extract does not identify \(H\) sufficiently to decide a single word problem, and that an associated Leary group family has a criterion depending on the base group and on whether \(S\) is recursive; it therefore describes the unqualified claim as unresolved. The background calls this dated literature triage derived from supplied Oberwolfach curation and requests independent verification of time-sensitive status and cited claims.
\item The record labels the statement \texttt{corrected\_verified} and reports that the official report was checked directly and that the clean statement preserves the complete target while restoring definitions and removing extraction artifacts. These are supplied source-verification assertions; no source visit or independent verification is claimed in this draft.
\end{itemize}
\end{document}
